\section{Project Proposal}

I propose to develop a one-dimensional Chebyshev spectral collocation solver for the McNabb-Foster hydrogen transport equations \cite{mcnabb1963}, which couple Fickian diffusion of mobile hydrogen with trapping and detrapping kinetics across multiple continuously distributed trap populations in metals. The governing system is
\begin{align}
\frac{\partial c_m}{\partial t} &= \nabla \cdot \left( D(T) \nabla c_m \right) + S(x,t) - \sum_i \frac{\partial c_{t,i}}{\partial t}, \label{eq:mobile} \\
\frac{\partial c_{t,i}}{\partial t} &= \nu_{m,i}(T) \, c_m \left( n_i - c_{t,i} \right) - \nu_{r,i}(T) \, c_{t,i}, \label{eq:trapped}
\end{align}
where $c_m$ is the mobile hydrogen concentration, $c_{t,i}$ is the concentration of hydrogen in trap $i$ with density $n_i$, $D(T) = D_0 \exp(-E_D / k_B T)$ is the Arrhenius diffusivity, $S(x,t)$ is a volumetric source term, and $\nu_{m,i}$ and $\nu_{r,i}$ are the temperature-dependent trapping and detrapping rates, respectively. Because all of these coefficients depend exponentially on temperature, the system is stiff and computationally demanding. The current state-of-the-art tool for these simulations is FESTIM \cite{delaporte2024}, a FEniCS-based finite element code that discretizes the spatial domain with continuous Galerkin linear elements (CG1) and advances in time with an implicit Euler scheme. While FESTIM is accurate and flexible, its low-order spatial discretization converges only algebraically under mesh refinement, and in practice this means that resolving the steep concentration gradients near material surfaces and boundary condition transitions requires very fine meshes. Chebyshev spectral methods, by contrast, achieve exponential convergence for smooth solutions and naturally cluster collocation nodes near domain boundaries \cite{boyd2001}, which is precisely where the sharpest gradients appear in hydrogen transport problems. This is why I believe spectral collocation is a strong candidate for accelerating these simulations.

The solver will target two representative experimental configurations. First, thermo-desorption spectroscopy (TDS), where a temperature ramp drives trapped hydrogen out of a previously implanted sample, and the goal is to reproduce the desorption flux spectrum more efficiently than a standard finite element approach. Second, multi-temperature gas-driven permeation (GDP), where sequential permeation experiments at increasing temperatures (e.g., 400\,$^\circ$C, 500\,$^\circ$C, 600\,$^\circ$C) with evacuations between each stage can cause FESTIM's finite element solver to diverge when the upstream pressure boundary condition is reintroduced at higher temperatures, because the exponentially larger diffusivity produces steep boundary layers that low-order elements struggle to capture. I will benchmark the Chebyshev collocation solver against both FESTIM's default CG1 discretization and higher-order finite element schemes (CG2, CG3) to isolate the advantage of spectral convergence over conventional $p$-refinement. Comparisons will include spatial convergence rates, wall-clock time, and degrees of freedom required for equivalent precision, with time integration handled by implicit BDF methods to manage the stiffness of the Arrhenius kinetics.
